\documentclass{article}
\usepackage{graphicx}

\usepackage[utf8]{inputenc}
\usepackage[T1]{fontenc}
\usepackage{needspace}
\usepackage{marginnote}
\renewcommand*{\marginfont}{\sffamily\footnotesize}

\usepackage{imakeidx}
\usepackage{hyperref}
\makeindex[intoc]

\usepackage{background}


\backgroundsetup{
scale=1,
angle=34,
opacity=0.4,
contents={\begin{tikzpicture}[remember picture,overlay]
 \path [left color = red!35,middle color = red!25, right color = white,rotate=1] (current page.south west)rectangle (current page.north east);   % Adjust the position of the logo.
\end{tikzpicture}\begin{tikzpicture}[remember picture,overlay]
 \path [left color = white,middle color = blue!10, right color = blue!20,rotate=13] (current page.south west)rectangle (current page.north east);   % Adjust the position of the logo.
\end{tikzpicture}\begin{tikzpicture}[remember picture,overlay]
 \path [left color = white,middle color = yellow!10, right color = white ,rotate=-13] (current page.south west)rectangle (current page.north east);   % Adjust the position of the logo.
\end{tikzpicture}}
}
\newenvironment{overview}[3][]{
  \needspace{10\baselineskip}
  \begin{center}
    { \renewcommand\textsuperscript[1]{}
      \phantomsection\addcontentsline{toc}{section}
      {\texorpdfstring{#2 (\emph{#3})}{#2 (#3)}}
    }
    {{\large\bfseries #2}\marginnote{#1}\par}
    \medskip
    {#3\par}
    \smallskip

    
  {%
  \bigskip
  \hrule
  \bigskip
  }
  
  \end{center}
}
  
\begin{document}
\includegraphics[scale=0.3]{mat-mn-navn-eng.eps} 


\overview{2D or not 2D: Dimensional and mechanical cues in cell migration.}{Cinzia Anita Maria Progida}{The movement of cells within the body, or cell migration, is fundamental for both physiological and
pathological processes including wound healing, immune response and cancer metastasis. Cell
motility involves the reorganization of the cytoskeleton and the movement of organelles, usually
triggered by external cues. Indeed, cells respond to biochemical or mechanical stimuli by activating
signaling pathways, re-organizing the cytoskeleton and generating forces. We study how subcellular
components and their intracellular dynamics affect the cellular response to extracellular stimuli such
as dimensional cues and physical confinement and the molecular mechanisms involved in
modulating cell migration.
The majority of the cell migration studies have been conducted using two-dimensional (2D) systems,
where adhesion formation and turnover represent important steps. However, cells require different
migratory strategies in 2D systems compared to the more physiological 3D systems. In addition,
some type of cells, such as immune cells, exhibit an adhesion-independent migration strategy in
their native environment.
I will present the different migration assays we use in order to characterize how subcellular
components influence the cellular response to dimensional cues and physical confinement. We
either embed cells in collagen gels or we use micro-fabricated channels where the cells are
completely confined. We also study how sub-cellular- level processes affect the forces exerted by the
cells on their environment using micropillars.
}



\overview{Membrane formation and compartmentalization in synthetic active matter}{Liesbeth M. C. Janssen}{

Active matter refers to systems whose constituent agents can move autonomously through the
consumption of energy. Such materials exhibit rich non-equilibrium dynamics and provide a framework
to understand the complex collective behavior seen in many living systems. In this talk, I will highlight
recent results of particle-resolved simulations of active rods that mimic cell-like properties such as
membrane formation and compartmentalization of particles. The crucial ingredients in our model are
the particles’ intrinsic self-propulsion, interparticle aligning interactions, and an inhomogeneous
motility field. We thus show that a minimal artificial active-matter system can self-organize into
structures reminiscent of biological patterns. Our predictions may be verified experimentally in e.g.
vibrated granular matter and other dry active systems with a spatially dependent self-propulsion speed.

\small{
Reference: J. Grauer, H. Löwen, and L.M.C. Janssen, Spontaneous membrane formation and self-
encapsulation of active rods in an inhomogeneous motility field, arXiv:1707.03405.}
}


\overview{
Mechanics of the cell interface studied by supported lipid bilayers
}{Margarita Staykova}{ The cell membrane undergoes complex morphological and surface area transformations while being confined to an underlying actin cortex, the membranes of neighboring cells or other extracellular structures. To understand the role of confinement in the membrane processes we adhere synthetic lipid bilayers to artificial substrates and subject them to perturbations that are common to the cell membrane - 1) substrate area changes, or 2) intake of extra lipids. Our results show that confined lipid bilayers regulate the arising changes in their lipid density either by the expulsion and absorption of lipid protrusions, such as tubes and vesicles, or by sliding over the substrate. Similar processes have been recently confirmed in living cells. We provide a theoretical framework that rationalizes the membrane behavior in terms of the membrane elasticity, and the adhesion and hydrodynamic interactions between the membrane and the substrate. 
}



\overview{The role of biomechanical modelling in discovering the brain’s lymphatic system}{Alexandra K. Diem}{How does the human brain eliminate waste products? This simple yet crucial question has
occupied medical research for decades, in particular with regards Alzheimer's disease,
whose onset is closely associated with a failure to remove the cerebral waste product
amyloid-$\beta$ (A$\beta$). Analytical and numerical modelling of the biomechanical processes
underlying waste removal from the brain can play a crucial role in evaluating hypotheses
and identifying its driving mechanisms. This talk describes the journey of attempting to
resolve some of the questions surrounding the removal of A$\beta$ in healthy individuals,
starting from the biomedical evidence, to the analytical and numerical methods to disprove
one of the most popular hypotheses and the derivation of new hypotheses. I will argue that
in order to fully resolve the onset and development of neurodegenerative diseases we
require fully interdisciplinary scientists working side-by- side with experimental researchers
and taking into account the multiple scales involved in the cerebral waste transport
processes in the human brain.}

\overview{}{Professor M\'ar M\'asson}{}


\overview{Biophysical approach of the mucociliary function: Mucus rheology and beating coordination}{ MASSIERA Gladys}{The mucociliary function of the bronchial epithelium ensures the continuous clearance of the respiratory system, which relies on two main elements: mucus and cilia beating coordination.
We perform here a rheological characterization of mucus samples extracted from ALI (Air-liquid interface) cultures of bronchial epithelium. Our approach combines  macro- and micro-rheology techniques with the aim of quantifying the mucus viscoelastic properties at different length scales (from the size of bronchial cilia up to the scale on which mucus is transported)
}




\overview{Hydration Effects Turn a Highly Stretched Polymer from an Entropic into an Energetic Spring}{Susanne Liese}{
Polyethylene glycol (PEG) is a structurally simple and nontoxic water-soluble polymer
that is widely used in medical and pharmaceutical applications as molecular linker
and spacer. In such applications, PEG’ s elastic response against conformational
deformations is key to its function. According to text-book knowledge, a polymer
reacts to the stretching of its end-to-end separation by a decrease in entropy that is
due to the reduction of available conformations, which is why polymers are commonly
called entropic springs. By a combination of single-molecule force spectroscopy
experiments with molecular dynamics simulations in explicit water, we show that
entropic hydration eff ects almost exactly compensate the chain conformational
entropy loss at high stretching. Our simulations reveal that this entropic
compensation is due to the stretching-induced release of water molecules that in the
relaxed state form double hydrogen bonds with PEG. As a consequence, the
stretching response of PEG is predominantly of energetic, not of entropic, origin at
high forces and caused by hydration eff ects, while PEG backbone deformations only
play a minor role. These fi ndings demonstrate the importance of hydration for the
mechanics of macromolecules and constitute a case example that sheds light on the
antagonistic interplay of conformational and hydration degrees of freedom.}


\overview{Transport of nutrients and clearance of waste products through the tiny
spaces surrounding brain cells}{Klas Pettersen}{Transport of nutrients and clearance of waste products are
prerequisites for healthy brain function. The brain lacks lymph vessels and must
rely on other mechanisms for clearance of waste products, including
amyloid $\beta$ that may form pathological aggregates if not effectively cleared. It is
still debated whether solutes are transported through the tiny spaces
surrounding brain cells, the interstitial space, by pressure-mediated bulk flow or
by diffusion. In a recently published article we simulated interstitial bulk flow
within 3D electron microscope reconstructions of hippocampal tissue [1]. We
found that the permeability is one to two orders of magnitudes lower than values
typically seen in the literature, arguing against bulk flow as the dominant
transport mechanism. Further, we showed that solutes of all sizes are more
easily transported through the interstitium by diffusion than by bulk flow. We
conclude that diffusion within the interstitial space combined with advection
along vessels is likely to substitute for the lymphatic drainage system in other
organs.


\small{[1] Holter KE, Kehlet B, Devor A, Sejnowski TJ, Dale AM, Omholt SW, Ottersen OP,
Nagelhus EA, Mardal K-A and Pettersen KH (2017). Interstitial solute transport
in 3D reconstructed neuropil occurs by diffusion rather than bulk flow.
Proceedings of the National Academy of Sciences, 6, 201706942.
http://doi.org/10.1073/pnas.1706942114}
}

\overview{Upside-Down and Inside-Out: The Biomechanics of Cell Sheet Folding}{Professor Raymond E. Goldstein}{The mucociliary function of the bronchial epithelium ensures the continuous clearance of the respiratory system, which relies on two main elements: mucus and cilia beating coordination.
We perform here a rheological characterization of mucus samples extracted from ALI (Air-liquid interface) cultures of bronchial epithelium. Our approach combines  macro- and micro-rheology techniques with the aim of quantifying the mucus viscoelastic properties at different length scales (from the size of bronchial cilia up to the scale on which mucus is transported).}


\overview{Physics of epithelial folding}{Guillaume Salbreux}{  
Three-dimensional deformations of epithelia play a fundamental role in tissue morphogenesis. The shape of an epithelium is determined by mechanical stresses acting within the tissue cells and from the outside environment. Here we introduce a three-dimensional vertex model which allows to represent the shape of a tissue in three dimensions by a set of vertices. In the model, the motion of vertices is set by apical, lateral and basal surface and line tensions, as well as intracellular pressures and external forces. Using this framework, we discuss how patterned force generation in an epithelium can drive biological tissue folding in fold formation in the Drosophila wing disc and in pancreatic tumour formation.}


\overview{Computational design of artificial creatures}{Gazzola, Mattia}{We introduce an inverse design approach based on minimal theoretical modeling, direct numerical simulations and artificial intelligence for the investigation of animal locomotion. We will mostly focus on aquatic and terrestrial limbless creatures and discuss the identification of optimal swimming gaits and morphologies, as well as the design of cyborg creatures.}


\overview{A cell-based framework for numerical modeling of electrical conduction in cardiac tissue}{Aslak Tveito}{In every heartbeat, an electrical wave traverses the entire heart muscle, and in every heart cell,
this wave sets off an action potential that increases the transmembrane potential of the cell.
This leads to the release of large amounts of calcium from internal storage structures of the cell.
Increased cytosolic calcium concentration, in turn, leads to contraction of the cells, and the well
synchronized contraction of about nine billion cells underlie the pumping function of the heart.
The electrical wave originating in the Sino-atrial node and spreading, at high speed, throughout
the heart muscle is of vital importance of every human and perturbations to the wave – referred
to as arrhythmias – can be life-threatening. Significant efforts are therefore invested in
understanding this wave and how dangerous perturbations can be avoided.
Over the last 60 years mathematical models have been used intensively to study electrical
conduction in the heart. The models involved are based on homogenization of the cardiac tissue
and are of reaction diffusion type where a parabolic equation is coupled to a large system of
ordinary differential equations defined in every point of the tissue. These models have been
very successful in providing a basic understanding of the properties of the excitation waves, and
have represented the level of accuracy that has been computationally feasible.
With increasing computing capacities, it has become clear that ever more realistic models can
be solved. Such models can increase insight into the astonishingly complex processes underlying
every heartbeat. In the present talk, we will discuss a cell-based framework for modeling the
electrical conduction system which avoids parts of the homogenization used in the classical
models. We will discuss the computational problems arising in the cell-based models and
provide some examples of simulations revealing new insight into electrical conduction through a
small number of cells.}



\overview{Modeling blood flow and mass transfers in cerebral microcirculation
}{Sylvie Lorthois}{After underlying the central role of cerebral microcirculation in brain physiology and in several brain pathologies, I will present the architecture of the microvascular cerebral network and show that it is the superposition of two types of structures: a mesh-like capillary structure, homogeneous over a cut-of length corresponding to the characteristic length of capillary vessels 50 $\mu$m), and fractal arborescent structures composed of arteries and veins.

Based on these results, I will present some of the approaches we develop for studying blood flow and/or or mass transfer at various scales, most of which are based on methodologies developed for the study of multiphase or reactive flows in porous media. Finally, I will present some perspectives related to the role of cerebral microcirculation in neurodegenerative diseases.}


\overview{Real time imaging of nanoparticles during tuberculosis infection in zebrafish embryos.}{Federico Fenaroli}{The use of nanotechnology is set to change the way we treat diseases in the decades to come. In particular, one of the areas which appear to be among the most promising is the treatment of Tuberculosis (TB), caused by the bacterium Mycobacterium tuberculosis. This optimism is based on a series of experiments that have tested nanoparticles (NP) enclosing antibiotics that were carried out in different mammalian models of TB, from mouse to monkey. However, in none of these studies was it possible to follow the localisation of bacteria and the NP inside the animals; the evaluation of treatment was based solely on the number of bacteria that survived after sacrificing the animals.  For this reason our group has introduced the use of the transparent zebrafish embryo and a fish model of TB as a powerful tool to visualize both bacteria and NP in detail and in real time in live animals. For this we use the fish TB pathogen Mycobacterium marinum (Mm), the causative agent of TB in fish, and NP made of different polymers such as poly lactide co-glycolide, or lipid-based liposomes (1). Routinely, the bacteria express red fluorescent protein whereas the NP carry a green fluorescent dye,

The zebrafish embryo possesses several unique characteristics that render it an ideal animal model when studying NP flow dynamics and their interactions with infected cells. The model vertebrate is easy to manipulate genetically and, as a consequence, several fish lines are now available that exhibit selectively labelled cells of the immune system. These features give a unique advantage to the zebrafish larvae for non-invasive and simultaneous observation of the pathogen, NPs and immune cells in a living vertebrate using fluorescence microscopy. The interactions among these players can also be observed using a technique we recently introduced in the zebrafish embryo, optical tweezers, which allows for in vivo manipulation of NPs (2). When studying TB in zebrafish, we observed that the injection of fluorescent liposomes into Mm-infected fish led to passage of NP from the blood across the endothelial cell boundary into the area of infection in a time dependent manner. Further quantitative analysis has shown the importance of the surface and NP size when measuring their accumulation in the diseased area. The ability to monitor NP, bacterial infection and different parameters related to flow of blood through blood vessels in vivo at unprecedented resolution opens up a powerful system for use of theoretical modelling in future studies.\\
1) Fenaroli, F. et al. : ACS Nano 2014, 8(7): 7014-7026.\\
2) Johansen, P. et al. : Nature communications 2016, 7:10974.
}


\overview{}{}{}


\overview{}{}{}

\overview{}{}{}

\overview{Rognes}{}{}


\overview{}{}{}







\end{document}